\begin{corrige}{4-12}

\begin{enumerate}

\item Soit $X$ de loi $\mathcal{N}(5,2)$. On pose: $$Z = \frac{X-5}{\sqrt{2}}$$ et $Z$ suit la loi normale centr\'ee r\'eduite. On calcule alors:
$$
P(4 \leq X \leq 6) = P ( \frac{4-5}{\sqrt{2}} \leq \frac{X-5}{\sqrt{2}} \leq \frac{6-5}{\sqrt{2}}) = 
P ( -0.707 \leq \frac{X-5}{\sqrt{2}} \leq 0.707 ) = 0.52,
$$

$$
P ( X  \geq 10 ) = P ( \frac{X-5}{\sqrt{2}} \geq \frac{10-5}{\sqrt{2}} ) = 2.03.10^{-4},
$$

$$
P(5.8 \leq X \leq 6.2) = P ( \frac{5.8-5}{\sqrt{2}} \leq \frac{X-5}{\sqrt{2}} \leq \frac{6.2-5}{\sqrt{2}}) = 8.8.10^{-2}.
$$


\item De même, pour $X$ de loi $\mathcal{N}(-10,0.1)$. On pose: $$Z = \frac{X+10}{\sqrt{0.1}}$$ et $Z$ suit la loi normale centr\'ee r\'eduite. Puis on calcule:
$$
P(-11.5 \leq X \leq -8.5) = P ( \frac{-11.5+10}{\sqrt{0.1}} \leq \frac{X+10}{\sqrt{0.1}} \leq \frac{-8.5+10}{\sqrt{0.1}}) = 0.9999979,
$$

$$
P ( X  \leq -8 ) = P ( \frac{X+10}{\sqrt{0.1}} \leq \frac{-8+10}{\sqrt{0.1}} ) = 1.000
$$

$$
P(-7.8 \leq X \leq -7.6) = P ( \frac{-7.8+10}{\sqrt{0.1}} \leq \frac{X+10}{\sqrt{0.1}} \leq \frac{-7.6+10}{\sqrt{0.1}}) = 1.7.10^{-12}.
$$


\end{enumerate}

\end{corrige}
